\section{Introduction}

Access control answers whether a principal may obtain data. Purpose limitation
adds a different question: whether the same data may affect this decision.
Financial institutions routinely hold information for investigation, support,
fraud, marketing, and regulatory purposes whose permissible uses do not
coincide. A confidential-computing boundary can protect bytes from an
infrastructure operator while still allowing the enclosed model to use the
wrong attribute. Removing every confidential attribute avoids that misuse but
can erase the legitimate value of protected computation.

This paper studies that tension through Purpose-Selective Bounded Execution
(PSBE). PSBE binds a declared purpose to exact data, workload/model identity,
capabilities, release rules, privacy budget, lifecycle state, and reconstructable
evidence. FinBoundBench asks whether the same synthetic confidential signal can
improve an approved financial task and cease to influence a distinct prohibited
task. The comparison includes full-data, prompt-only, a strong deterministic
prefilter, and layered PSBE conditions. In a three-phase OpenRouter diagnostic,
three of five models pass admission, 83\% of position diagnostic invocations
pass, and 89\% of confirmatory invocations pass. Release denials are
concentrated in one model, suggesting model-specific rather than system-level
failure. The intended contributions are a formal definition, paired benchmark,
cross-layer attack suite, and empirical analysis; the full confirmatory claim
family requires the registered 200-pair run.
